\section{Introduction}
\label{sec:intro}

% TODO: motivation. Compression is needed to deploy medical VLMs on limited hardware.
% Papers report average accuracy; clinically important failures (small lesions, rare
% findings, answering without looking at the image, overconfidence, hallucinated
% findings) can get worse while average accuracy holds.
\todo{motivation: why average accuracy hides compression harm}

% TODO: research questions.
% RQ1: At a matched measured budget, do compression families with similar average
%      accuracy differ in small-finding recall, rare-finding recall, image reliance,
%      calibration and hallucination?
% RQ2: Which compression family (quantization, distillation, generic pruning,
%      layer pruning) damages which property?
% RQ3: Does a clinically weighted pruning criterion preserve small findings and
%      image reliance better than generic criteria at the same budget?
\todo{RQ1--RQ3}

% TODO: contributions.
% C1 Audit: controlled comparison of E1-E9 at matched measured budgets on hidden
%    metrics, not only accuracy.
% C2 Clinical pruning criterion (E7) with pre-registered finding weights.
% C3 Open evaluation harness (stress tests, mask audit, efficiency checks).
\todo{contributions C1--C3}

\begin{figure}[t]
  \centering
  \fbox{\parbox{0.9\linewidth}{\centering\vspace{3em}
    Placeholder: headline accuracy vs.\ small-finding recall for \expid{E2},
    \expid{E3}, \expid{E5}, \expid{E6}, \expid{E7} at \expid{B2}.\vspace{3em}}}
  \caption{Teaser: compression methods with similar average accuracy differ in
  small-finding recall. Fed by \expid{E2}, \expid{E3}, \expid{E5}--\expid{E7} at
  \expid{B2}.}
  \label{fig:teaser}
\end{figure}
